%=============================================================================
\section{Conclusion}
\label{sec:conclusion}
%=============================================================================

We have applied persistent homology to co-occurrence networks derived from
3,409 fecal samples in the American Gut Project, using a paired bootstrap
design with explicit FDR correction and covariate matching to assess the
robustness of group differences in topological structure.

The primary finding is clear: inflammatory bowel disease is associated with
substantially lower topological complexity in microbial co-occurrence networks,
with large and fully robust effect sizes ($d = 0.75$--$2.38$) across all six
\Hone\ features, all three taxon-set sizes tested ($N \in \{50, 80, 120\}$),
and both full and matched subsamples.
This constitutes, to our knowledge, the first demonstration that \ibd\
dysbiosis has a detectable topological signature in the loop structure of
co-occurrence networks.

A secondary finding is that plant-based diet is associated with longer-lived
co-occurrence loops, suggesting a structural reorganization of the microbial
interaction landscape rather than simple enrichment or depletion.
This finding is partially robust to covariate matching and stable across two
of three taxon-set sizes.

We make no claim about antibiotic recency based on these data; the signal is
weak and inconsistent.

The paired bootstrap design introduced here---drawing matched subsamples
repeatedly, computing topological features per iteration, and testing the
distribution of paired differences against a sign-flip null---is directly
reusable for any group comparison in microbiome TDA.
Code and data are publicly available at
\href{https://github.com/adam-jfkhs/Microbiome-TDA}{\texttt{github.com/adam-jfkhs/Microbiome-TDA}}.

\paragraph{Future directions.}
Immediate priorities are: (1) correlation of topological features with
conventional alpha diversity and beta diversity metrics to establish
the informational overlap and added value; (2) taxon-level attribution of
loop structure to identify which co-occurrence pairs form the persistent
\Hone\ features; (3) replication in an independent IBD cohort; and (4)
longitudinal analysis to assess whether topological simplification precedes
or follows clinical relapse.
